% Source PDF: source/普通高考/2013/2013江西理.pdf, page 1, question 7.
% pdfimages -list found no embedded images; the source program flowchart is vector artwork.
% Node -> directed-edge list: 开始 -> i=1,S=0 -> i=i+1 -> i 是奇数;
% 是 -> blank process box; 否 -> S=2*i+1; both feed S<10;
% S<10 --是--> loop back to the i=i+1 trunk; S<10 --否--> 输出 i -> 结束.
% The blank process box is intentionally left blank because the source box is blank.
% All borders and directed connectors are solid; node text is centered with local parindent=0pt.

\begin{tikzpicture}[
    >=Stealth,
    x=1cm,
    y=1cm,
    line width=0.45pt,
    line cap=round,
    line join=round,
    font=\normalsize,
    flowtext/.style={anchor=center, align=center,
        execute at begin node={\setlength{\parindent}{0pt}}},
    flowbox/.style={flowtext, draw, minimum height=0.68cm,
        inner xsep=4pt, inner ysep=2pt},
    terminal/.style={flowbox, rounded corners=2pt, minimum width=1.35cm},
    process/.style={flowbox, minimum width=2.50cm},
    decision/.style={flowbox, diamond, aspect=2.8, inner sep=0pt,
        minimum width=4.15cm, minimum height=1.05cm},
    io/.style={flowbox, trapezium, trapezium left angle=72, trapezium right angle=108,
        minimum width=1.80cm},
    branchlabel/.style={flowtext, inner sep=0pt},
]
    \node[terminal] (start) at (0,9.15) {开始};
    \node[process, minimum width=2.75cm] (init) at (0,7.85) {$i=1,S=0$};
    \node[process, minimum width=2.15cm] (inc) at (0,6.45) {$i=i+1$};
    \node[decision] (odd) at (0,5.05) {$i$ 是奇数};
    \node[process, minimum width=1.55cm] (blank) at (0,3.65) {};
    \node[process, minimum width=2.70cm] (oddcase) at (3.60,5.05) {$S=2*i+1$};
    \node[decision, minimum width=3.65cm] (test) at (0,2.10) {$S<10$};
    \node[io, minimum width=1.85cm] (output) at (0,0.65) {输出 $i$};
    \node[terminal] (finish) at (0,-0.70) {结束};

    \draw[->] (start.south) -- (init.north);
    \draw[->] (init.south) -- (inc.north);
    \draw[->] (inc.south) -- (odd.north);

    % Odd branch goes down through the blank source box; even branch exits
    % right to its assignment, then both enter the S<10 test separately.
    \draw[->] (odd.south) -- node[branchlabel, left=5pt,pos=0.18] {是} (blank.north);
    \draw[->] (odd.east) -- node[branchlabel, above=4pt,pos=0.18] {否} (oddcase.west);
    \draw[->] (blank.south) -- (test.north);
    \draw[->] (oddcase.south) -- (3.60,3.05) -- (0,3.05);

    % S<10 false exits to output; the true branch loops around the left and
    % returns with a right-pointing arrow above the independent trunk arrow.
    \draw[->] (test.south) -- node[branchlabel, left=2pt] {否} (output.north);
    \draw[->] (output.south) -- (finish.north);
    \draw (test.west) -- (-2.85,2.10) -- (-2.85,7.15);
    \draw[->] (-2.85,7.15) -- (0,7.15);
    \node[branchlabel] at (-1.98,2.38) {是};

    % Keep a tight bounding box around the outer loop and terminals.
    \path[use as bounding box] (-3.05,-1.05) rectangle (5.00,9.45);
\end{tikzpicture}
